\section{Banco de questões — Infraestrutura de TI}

\subsection*{N1 — Fundamento competitivo}

\begin{questao}{\nivel{N1} 10.01 — Cebraspe/Certo ou Errado}
Em uma rede IPv4 \texttt{/27}, o bloco contém 32 endereços; em uso convencional, 30 podem ser atribuídos a hosts, pois um endereço identifica a rede e outro o broadcast.
\end{questao}

\begin{questao}{\nivel{N1} 10.02 — Cebraspe/Certo ou Errado}
O fato de o UDP não fornecer retransmissão e ordenação nativamente implica que qualquer aplicação baseada em UDP seja, necessariamente, incapaz de oferecer entrega confiável.
\end{questao}

\begin{questao}{\nivel{N1} 10.03 — Cebraspe/Certo ou Errado}
Uma VLAN separa domínios de broadcast em camada 2, mas não impede, por si só, comunicação entre redes quando existe roteamento inter-VLAN com política permissiva.
\end{questao}

\begin{questao}{\nivel{N1} 10.04 — Cebraspe/Certo ou Errado}
Em um switch Ethernet, o endereço MAC de origem recebido em um quadro é usado para aprendizado da tabela de encaminhamento, ao passo que o endereço MAC de destino é usado na decisão de encaminhamento.
\end{questao}

\begin{questao}{\nivel{N1} 10.05 — Cebraspe/Certo ou Errado}
SNMPv3 pode fornecer autenticação e privacidade, mas sua mera adoção não garante gerenciamento seguro se credenciais, perfis e acesso à rede de gerenciamento forem mal configurados.
\end{questao}

\begin{questao}{\nivel{N1} 10.06 — Cebraspe/Certo ou Errado}
LDAP e Active Directory são termos equivalentes, pois ambos designam um serviço de diretório completo e independente de outros protocolos.
\end{questao}

\begin{questao}{\nivel{N1} 10.07 — Cebraspe/Certo ou Errado}
RAID 6 tolera, em condições normais do arranjo, a falha de até dois discos, mas isso não o transforma em substituto de backup.
\end{questao}

\begin{questao}{\nivel{N1} 10.08 — Cebraspe/Certo ou Errado}
Uma snapshot armazenada no mesmo storage da máquina virtual constitui cópia independente suficiente contra perda total do storage original.
\end{questao}

\begin{questao}{\nivel{N1} 10.09 — Cebraspe/Certo ou Errado}
Em um sistema paginado, thrashing caracteriza situação em que o tempo gasto com faltas e substituição de páginas cresce a ponto de reduzir significativamente o trabalho útil da CPU.
\end{questao}

\begin{questao}{\nivel{N1} 10.10 — Cebraspe/Certo ou Errado}
Duas máquinas virtuais executadas no mesmo host físico eliminam o ponto único de falha do hardware, pois cada VM constitui um domínio de falha independente.
\end{questao}

\subsection*{N2 — Aplicação}

\begin{questao}{\nivel{N2} 10.11 — FCC/subnetting}
Um segmento precisa suportar 50 hosts IPv4 em uma única sub-rede, sem considerar endereçamento especial além de rede e broadcast. O menor prefixo CIDR adequado é:
\begin{enumerate}[label=\Alph*)]
\item /28
\item /27
\item /26
\item /25
\item /24
\end{enumerate}
\end{questao}

\begin{questao}{\nivel{N2} 10.12 — FGV/roteamento}
A tabela de rotas de um roteador contém \texttt{10.0.0.0/8}, \texttt{10.20.0.0/16} e \texttt{10.20.30.0/24}. Um pacote destinado a \texttt{10.20.30.55} será encaminhado, em regra, pela rota:
\begin{enumerate}[label=\Alph*)]
\item \texttt{10.0.0.0/8}, por ser a mais antiga;
\item \texttt{10.20.0.0/16}, por possuir métrica intermediária;
\item \texttt{10.20.30.0/24}, por ser o prefixo mais específico;
\item default, porque existem rotas sobrepostas;
\item aleatoriamente entre as três.
\end{enumerate}
\end{questao}

\begin{questao}{\nivel{N2} 10.13 — Cebraspe/Certo ou Errado}
Um trunk 802.1Q permite transportar múltiplas VLANs entre switches, mas não substitui o roteamento necessário para comunicação entre VLANs distintas.
\end{questao}

\begin{questao}{\nivel{N2} 10.14 — FGV/802.1X}
Uma instituição deseja liberar portas cabeadas e acesso Wi-Fi apenas após autenticação centralizada. A arquitetura mais aderente é:
\begin{enumerate}[label=\Alph*)]
\item 802.1X com supplicant, autenticador e backend RADIUS/EAP;
\item STP com DNS recursivo;
\item SNMP com community string pública;
\item LACP com NAT;
\item MPLS com DHCP relay.
\end{enumerate}
\end{questao}

\begin{questao}{\nivel{N2} 10.15 — Cebraspe/Certo ou Errado}
O bloqueio indiscriminado de ICMP pode prejudicar funcionalidades legítimas de rede, como diagnóstico e mecanismos relacionados a descoberta de MTU de caminho.
\end{questao}

\begin{questao}{\nivel{N2} 10.16 — FCC/processos}
Um processo em execução solicita leitura de disco e não pode prosseguir até a conclusão do I/O. O estado mais coerente, imediatamente após a solicitação, é:
\begin{enumerate}[label=\Alph*)]
\item pronto;
\item bloqueado/esperando;
\item terminado;
\item zombie;
\item executando com prioridade máxima.
\end{enumerate}
\end{questao}

\begin{questao}{\nivel{N2} 10.17 — FGV/deadlock}
Dois processos mantêm recursos distintos e cada um espera indefinidamente pelo recurso mantido pelo outro. O cenário evidencia diretamente:
\begin{enumerate}[label=\Alph*)]
\item thrashing;
\item starvation sem espera circular;
\item uma cadeia de espera circular compatível com deadlock;
\item fragmentação externa;
\item falha de paginação.
\end{enumerate}
\end{questao}

\begin{questao}{\nivel{N2} 10.18 — Cebraspe/Certo ou Errado}
Em Linux, conceder \texttt{chmod 777} para resolver um erro de acesso pode eliminar a causa imediata, mas cria risco de privilégio excessivo e não deve ser tratado como solução segura por padrão.
\end{questao}

\begin{questao}{\nivel{N2} 10.19 — FGV/AD}
Em um domínio Active Directory, usuários conseguem autenticar-se localmente, mas falham ao localizar controladores de domínio após alteração incorreta do DNS. Isso é coerente porque:
\begin{enumerate}[label=\Alph*)]
\item o AD depende de DNS para localização de serviços e controladores;
\item LDAP substitui integralmente o DNS;
\item Kerberos não utiliza nomes de serviço;
\item GPO dispensa resolução de nomes;
\item o Global Catalog funciona sem qualquer infraestrutura DNS.
\end{enumerate}
\end{questao}

\begin{questao}{\nivel{N2} 10.20 — FCC/RAID}
Um arranjo RAID 5 utiliza quatro discos de 6 TB. Desprezando overhead e diferenças de unidade, a capacidade útil aproximada é:
\begin{enumerate}[label=\Alph*)]
\item 6 TB
\item 12 TB
\item 18 TB
\item 24 TB
\item 30 TB
\end{enumerate}
\end{questao}

\begin{questao}{\nivel{N2} 10.21 — Cebraspe/Certo ou Errado}
Em RAID 10 com quatro discos, a tolerância a duas falhas simultâneas depende de quais discos falham; perder os dois membros do mesmo espelho pode causar perda do conjunto.
\end{questao}

\begin{questao}{\nivel{N2} 10.22 — FGV/NAS e SAN}
Uma aplicação exige volume em bloco apresentado ao servidor como LUN, enquanto outra precisa de compartilhamento de arquivos SMB para usuários. As tecnologias mais aderentes são, respectivamente:
\begin{enumerate}[label=\Alph*)]
\item NAS e SAN;
\item SAN e NAS;
\item CDN e SAN;
\item SAN e DNS;
\item NAS e DHCP.
\end{enumerate}
\end{questao}

\begin{questao}{\nivel{N2} 10.23 — FCC/backup}
Uma estratégia utiliza backup full aos domingos e incrementais de segunda a sábado. Para restaurar o estado de sexta-feira, normalmente será necessário:
\begin{enumerate}[label=\Alph*)]
\item apenas o incremental de sexta;
\item o full de domingo e todos os incrementais de segunda a sexta;
\item o full de domingo e somente o incremental de segunda;
\item todos os fulls do mês;
\item apenas o último snapshot de memória.
\end{enumerate}
\end{questao}

\begin{questao}{\nivel{N2} 10.24 — Cebraspe/Certo ou Errado}
Um RPO de 15 minutos pode ser incompatível com a realização de apenas um backup completo diário, na ausência de mecanismos adicionais de replicação ou captura contínua de alterações.
\end{questao}

\begin{questao}{\nivel{N2} 10.25 — FGV/containers}
Uma aplicação mantém estado de sessão exclusivamente na memória local de cada contêiner e passa a ser escalada horizontalmente. A medida arquitetural mais adequada é:
\begin{enumerate}[label=\Alph*)]
\item fixar o tráfego permanentemente a um único contêiner;
\item externalizar/replicar o estado ou adotar desenho stateless compatível;
\item desabilitar health checks;
\item substituir TCP por UDP;
\item armazenar sessão em camada temporária da imagem.
\end{enumerate}
\end{questao}

\subsection*{N3 — Integração de concurso}

\begin{questao}{\nivel{N3} 10.26 — FGV/assertivas de redes}
Considere:

I. STP/RSTP busca impedir loops de camada 2.

II. LACP pode agregar enlaces físicos em um canal lógico.

III. A existência de LACP implica que um único fluxo TCP utilizará sempre a soma integral da largura de banda de todos os membros.

Assinale a alternativa correta.
\begin{enumerate}[label=\Alph*)]
\item Apenas I.
\item Apenas II.
\item Apenas I e II.
\item Apenas II e III.
\item I, II e III.
\end{enumerate}
\end{questao}

\begin{questao}{\nivel{N3} 10.27 — Cebraspe/Certo ou Errado}
Uma rede segmentada em VLANs continua vulnerável a movimento lateral se o firewall ou roteador inter-VLAN permitir tráfego amplo entre as zonas, pois a segmentação lógica sem política restritiva não garante isolamento efetivo.
\end{questao}

\begin{questao}{\nivel{N3} 10.28 — FCC/desempenho TCP}
Um link de 1 Gbit/s entre dois datacenters possui RTT elevado e baixa perda. Uma única transferência TCP atinge apenas fração da banda disponível. A explicação mais adequada é:
\begin{enumerate}[label=\Alph*)]
\item banda nominal é o único fator relevante, logo o comportamento é impossível;
\item janela TCP, RTT, congestionamento e produto largura de banda-atraso podem limitar o throughput de um fluxo;
\item UDP resolveria automaticamente qualquer problema de confiabilidade;
\item VLAN reduz throughput de longa distância por definição;
\item DNS é o principal limitador de transferência contínua após a conexão.
\end{enumerate}
\end{questao}

\begin{questao}{\nivel{N3} 10.29 — Cebraspe/Certo ou Errado}
O uso de MPLS VPN por uma operadora fornece isolamento lógico de tráfego, mas não deve ser confundido automaticamente com criptografia fim a fim do conteúdo transportado.
\end{questao}

\begin{questao}{\nivel{N3} 10.30 — FGV/SNMP}
Uma rede de gerenciamento usa SNMPv3, mas todos os equipamentos compartilham o mesmo usuário, a mesma chave e aceitam acesso de qualquer VLAN. A melhor avaliação é:
\begin{enumerate}[label=\Alph*)]
\item SNMPv3 torna o desenho seguro independentemente da configuração;
\item o protocolo oferece controles melhores, mas a configuração ainda viola segregação e mínimo privilégio;
\item deve-se voltar a SNMPv2c para reduzir risco;
\item MIB e SMI substituem autenticação;
\item community strings tornam-se mais seguras em SNMPv3.
\end{enumerate}
\end{questao}

\begin{questao}{\nivel{N3} 10.31 — FCC/memória virtual}
Um servidor apresenta alta utilização de disco, taxa elevada de page faults e CPU frequentemente ociosa esperando I/O. Aumentar apenas a frequência da CPU tende a produzir pouco efeito porque o gargalo principal pode ser:
\begin{enumerate}[label=\Alph*)]
\item excesso de memória disponível;
\item thrashing por pressão de memória/working set;
\item tabela MAC cheia;
\item DNS recursivo;
\item RAID 1.
\end{enumerate}
\end{questao}

\begin{questao}{\nivel{N3} 10.32 — Cebraspe/Certo ou Errado}
MV e contêiner oferecem isolamento por mecanismos distintos; em geral, uma VM inclui sistema operacional convidado sobre hipervisor, enquanto contêineres compartilham o kernel do host.
\end{questao}

\begin{questao}{\nivel{N3} 10.33 — FGV/HA}
Dois servidores de aplicação estão em hosts físicos distintos, mas ambos dependem do mesmo storage sem controladora redundante. A arquitetura possui:
\begin{enumerate}[label=\Alph*)]
\item alta disponibilidade completa, pois o compute é redundante;
\item ponto único de falha no storage compartilhado;
\item redundância 2N em todas as camadas;
\item tolerância a desastre regional;
\item independência total entre os servidores.
\end{enumerate}
\end{questao}

\begin{questao}{\nivel{N3} 10.34 — Cebraspe/Certo ou Errado}
A disponibilidade de um serviço não pode ser inferida apenas pela redundância de servidores, porque dependências comuns de DNS, storage, banco, energia ou telecomunicações podem permanecer como pontos únicos de falha.
\end{questao}

\begin{questao}{\nivel{N3} 10.35 — FCC/RPO-RTO}
Um sistema possui RPO de 30 minutos e RTO de 1 hora. O restore test mais recente recuperou dados de 10 minutos antes da falha, mas o serviço demorou 3 horas para voltar. Nesse teste:
\begin{enumerate}[label=\Alph*)]
\item RPO e RTO foram cumpridos;
\item somente o RPO foi cumprido;
\item somente o RTO foi cumprido;
\item nenhum foi cumprido;
\item não é possível avaliar nenhum dos dois.
\end{enumerate}
\end{questao}

\begin{questao}{\nivel{N3} 10.36 — FGV/backup}
Um banco usa RAID 6, snapshots locais a cada hora e backup externo semanal. A administração afirma que o risco de ransomware está totalmente coberto. A melhor avaliação é:
\begin{enumerate}[label=\Alph*)]
\item correta, porque RAID 6 protege contra ransomware;
\item correta, porque snapshots locais são imunes à exclusão;
\item inadequada; é preciso avaliar independência, imutabilidade, credenciais e janela de perda do backup externo;
\item correta se o banco for ACID;
\item correta se o storage usar SSD.
\end{enumerate}
\end{questao}

\begin{questao}{\nivel{N3} 10.37 — Cebraspe/Certo ou Errado}
A deduplicação pode reduzir consumo de armazenamento, mas aumenta dependência de metadados e não substitui a existência de cópias em domínios de falha independentes.
\end{questao}

\begin{questao}{\nivel{N3} 10.38 — FGV/Kerberos}
Uma conta administrativa utiliza a mesma identidade para e-mail, navegação e administração de controladores de domínio. Ainda que Kerberos funcione corretamente, o desenho é inadequado porque:
\begin{enumerate}[label=\Alph*)]
\item a mesma conta aumenta superfície de exposição de credencial privilegiada;
\item Kerberos proíbe e-mail;
\item LDAP deixa de funcionar;
\item GPO exige uma conta por aplicação;
\item DNS não suporta contas administrativas.
\end{enumerate}
\end{questao}

\begin{questao}{\nivel{N3} 10.39 — Cebraspe/Certo ou Errado}
Uma liveness probe que dependa de serviço externo pode provocar reinícios desnecessários de instâncias saudáveis durante indisponibilidade da dependência, amplificando a falha.
\end{questao}

\begin{questao}{\nivel{N3} 10.40 — FCC/capacidade}
A CPU média de um servidor é 22\%, mas o p99 de CPU chega a 100\% durante fechamento mensal, período em que o SLA é violado. A conclusão mais adequada é:
\begin{enumerate}[label=\Alph*)]
\item a média prova excesso de capacidade;
\item planejamento deve considerar picos, filas, percentis e janela crítica, não apenas média global;
\item CPU não influencia SLA;
\item o servidor está subutilizado em qualquer momento;
\item a única solução possível é scale-up.
\end{enumerate}
\end{questao}

\subsection*{N4 — Auditoria / avançado}

\begin{questao}{\nivel{N4} 10.41 — Cebraspe/caso integrado}
Em auditoria de continuidade, a equipe encontra: RAID 6; snapshots no mesmo storage; backup externo mensal; RPO declarado de 15 minutos; e nenhum restore test nos últimos 18 meses. Julgue o item:

A presença de RAID 6 e snapshots frequentes é suficiente para demonstrar o atendimento ao RPO, ainda que o storage seja perdido integralmente.
\end{questao}

\begin{questao}{\nivel{N4} 10.42 — FGV/rede corporativa}
Usuários, servidores e administração estão em VLANs distintas, porém a política do firewall interno é \texttt{any-any} entre todas as zonas. Em um teste de comprometimento de estação de usuário, o atacante alcança RDP e SMB dos servidores. O achado principal é:
\begin{enumerate}[label=\Alph*)]
\item inexistência de VLANs;
\item segmentação nominal sem política efetiva de contenção, permitindo movimento lateral;
\item falha exclusiva de DNS;
\item ausência de STP;
\item necessidade de trocar Ethernet por MPLS.
\end{enumerate}
\end{questao}

\begin{questao}{\nivel{N4} 10.43 — FCC/AD e logs}
Administradores de domínio usam contas compartilhadas e conseguem excluir logs locais dos controladores. Qual conjunto de controles é mais adequado?
\begin{enumerate}[label=\Alph*)]
\item contas nominativas, separação de privilégio, MFA/PAM e envio de logs para repositório protegido;
\item aumentar apenas o tamanho das senhas compartilhadas;
\item manter contas compartilhadas e reduzir retenção de logs;
\item desabilitar auditoria para melhorar desempenho;
\item migrar os controladores para RAID 0.
\end{enumerate}
\end{questao}

\begin{questao}{\nivel{N4} 10.44 — Cebraspe/Certo ou Errado}
Dois datacenters podem continuar compartilhando uma causa comum de falha se ambos dependem do mesmo provedor e do mesmo ponto físico de entrada de fibra, de modo que redundância geográfica aparente não garante diversidade de conectividade.
\end{questao}

\begin{questao}{\nivel{N4} 10.45 — FGV/containers}
Um cluster de contêineres aceita imagens sem assinatura, usa tag \texttt{latest} em produção e permite execução privilegiada por padrão. O risco combinado mais relevante é:
\begin{enumerate}[label=\Alph*)]
\item supply chain não verificável, baixa reprodutibilidade e privilégio excessivo;
\item apenas falta de espaço em disco;
\item somente indisponibilidade de DNS;
\item ausência de VLAN;
\item perda de capacidade RAID.
\end{enumerate}
\end{questao}

\begin{questao}{\nivel{N4} 10.46 — Cebraspe/Certo ou Errado}
Se todos os logs de infraestrutura permanecem apenas nos próprios servidores e podem ser apagados pelos mesmos administradores cujas ações são auditadas, a trilha de auditoria apresenta risco de integridade e accountability.
\end{questao}

\begin{questao}{\nivel{N4} 10.47 — FCC/HA}
Um serviço está em duas VMs no mesmo host, com banco em storage local do host e balanceador externo. A gestão classifica a arquitetura como altamente disponível. A avaliação correta é:
\begin{enumerate}[label=\Alph*)]
\item a classificação é adequada, pois existem duas VMs;
\item compute e dados compartilham o mesmo domínio de falha físico, logo a redundância é insuficiente;
\item o balanceador externo elimina qualquer SPOF interno;
\item o uso de virtualização substitui DR;
\item a única deficiência é o uso de storage local.
\end{enumerate}
\end{questao}

\begin{questao}{\nivel{N4} 10.48 — FGV/evidência de backup}
O software de backup registra ``job successful'' todas as noites. O auditor deseja concluir sobre recuperabilidade. A evidência mais apropriada é:
\begin{enumerate}[label=\Alph*)]
\item apenas o status verde do job;
\item teste documentado de restauração com validação de integridade e tempo obtido;
\item tamanho do arquivo de log;
\item capacidade bruta do storage;
\item existência de RAID no repositório.
\end{enumerate}
\end{questao}

\begin{questao}{\nivel{N4} 10.49 — Cebraspe/Certo ou Errado}
Um hash calculado sobre uma imagem forense e preservado durante a análise ajuda a demonstrar que a cópia não foi alterada, mas não prova, por si só, que a fonte original continha informação verdadeira.
\end{questao}

\begin{questao}{\nivel{N4} 10.50 — FGV/capacidade e risco}
Uma aplicação crítica possui p95 de 800 ms, p99 de 11 s e média de 320 ms. O SLA considera apenas média inferior a 1 s. Usuários relatam travamentos em operações de fechamento. A recomendação de auditoria mais adequada é:
\begin{enumerate}[label=\Alph*)]
\item manter o SLA, pois a média está cumprida;
\item revisar o indicador para incluir percentis e jornadas críticas, correlacionando latência, filas e saturação;
\item medir apenas disponibilidade;
\item remover monitoramento para reduzir overhead;
\item aumentar CPU sem diagnóstico adicional.
\end{enumerate}
\end{questao}
